\section{Overview}

We face three key research challenges in realizing hybrid simulation of ML frameworks: (1) supporting unmodified ML frameworks, (2) ensuring correct and efficient time synchronization, and (3) achieving scalability.
First, modern ML frameworks have complex software and hardware dependencies. They assume the presence of GPUs, NVLinks, and RDMA networks, and they rely on a wide range of libraries. Constructing an execution environment in which an ML framework can run unmodified is therefore nontrivial.
Second, integrating event-driven network simulation with ML framework execution presents a fundamental mismatch. Network simulators advance time in discrete time steps, while ML frameworks execute with continuous, real-time progression. Reconciling these two timing models is challenging.
Finally, ML frameworks are resource-intensive, consuming substantial compute, memory, and communication bandwidth. While we would like to execute them as-is, we must avoid the prohibitive resource costs of training a full-scale model on a real cluster.

\begin{figure}[t]
    \centering
    \includegraphics[width=.48\textwidth]{figures/overview.pdf}
    \caption{\sys architecture. Components in light green are unmodified, components in green are minimally modified, and components in blue are constructed by \sys.}
    \vspace{-2mm}
    \label{fig:overview}
\end{figure}


To support unmodified ML frameworks, \sys runs an ML framework in a realistic containerized environment and interacts with this real system for simulation. \autoref{fig:overview} shows \sys's architecture. Each container in the environment acts as a GPU server.
Each container runs Python interpreters that execute ML framework code.
There are two operations that require interaction between the containers and the \sys simulator: (1) GPU computation and (2) communication. When a computation kernel is invoked, \sys uses a single GPU to profile the kernel's execution time. For communication, \sys utilizes an event-driven flow-level network simulator to estimate completion time.
\sys only needs a single GPU because it profiles the performance once for each (computation kernel, tensor shapes) combination. This is sufficient because computation kernel performance is usually independent of the tensor values.\footnote{There are some exceptions. One such example is sorting where conditional branches could be chosen based on comparisons. We discuss this point in \autoref{sec:discussion}.}
In this way, an application cannot distinguish whether it is running on \sys or a physical GPU cluster as long as its control flow does not depend on tensor values (which would be junk values in \sys). The time of each rank in every container will be maintained by \sys using standard discrete event simulation, and the application can read this time to calculate its performance. A naive implementation of the above approach would still require modifying the ML framework source code. For example, one might need to change the performance timer implementation used in framework logging. To support ML frameworks out of the box, we instead leverage Python’s runtime patching to dynamically redirect underlying dependencies of an ML framework.

\begin{figure*}[t]
    \centering
    \includegraphics[width=.90\textwidth]{figures/workflow.pdf}
    \caption{An example workflow of \sys with two ranks. ML system places computation and communications on different CUDA streams for flexibility, and use CUDA events to manage the synchronization between them. \sys needs to correctly handle these synchronizations to achieve accurate simulation.}
    \vspace{-3mm}
    \label{fig:workflow}
\end{figure*}

To appreciate the time synchronization challenge, let's consider the following \textit{past events} scenario.
Traditional event-driven simulators require static workloads. For example, in a typical event-driven simulation workflow, all the events in the workload are pre-loaded into the event queue of the simulator. The simulator processes these events in chronological order, updating the simulation state and queuing new events as necessary. However, with real systems, events are generated dynamically and injected into the simulator, which can lead to past event scenarios: the simulator receives an event generated by the real system after the simulator has already advanced to the next event, which has a timestamp later than the real system's event.
One solution is to keep the event-driven network simulator tightly synchronized with the containers. In hardware simulators, for example, it is common to use a tiny time quantum and synchronize all components at every quantum~\cite{reinhardt1993wwt,mukherjee2000wwt}. However, it introduces significant overhead and defeats the purpose of event-driven simulation in modern network simulators.
\sys uses loose time synchronization between the event-driven simulator and the real system execution to address this issue and enables fast simulation. To ensure simulation accuracy, we implement an event-driven network simulator capable of time traveling to rollback the states of all flows when events occur in the past.


We apply two techniques to scale \sys. In LLM training, each GPU server may load training data and model weights, which consumes significant host memory. Our evaluation (\autoref{sec:eval-scalability}) shows that a physical server with 256GB host memory can only simulate 9 GPUs for the training of Llama2 7B model without our technique. We introduce a novel memory sharing mechanism for different containers running on the same host, which significantly reduces the memory footprint. 
Second, \sys measures only the actual CPU time consumed by each process, rather than the wall clock time. This choice preserves simulation accuracy in the presence of CPU core contention, where wall-clock measurements would otherwise overestimate execution time due to frequent context switches between the containers.
